\pagetopsection{ Witchcraft}


\label{sec:witchcraft}

\begin{multicols*}{2}

\paragraph{Patrons} are spirits, not gods. 

They don't have an overwhelming power over mortal realm, but rather significant control over their particular domain.
By virtue of being a spirit their grip on the material is indirect, that's why they need you. 

You lack power, but your actions can impact you surroundings, thus being a witch is a mutually beneficial agreement for you and for your patron.

\paragraph{Casting a spell} as a witch is a manifestation of favours granted to them by their patrons.
Favour is spent after a spell is complete and is regained when a witch indulges the passions of their patrons.

Cantrips don't cost favour, but stop working when a witch has none.

Each time after witch casts a spell in combat, they need to pass \texttt{Spirit Save} or become \texttt{Fatigued}. \textbf{No spell} can be cast until \texttt{Fatigue} is cleared.

\paragraph{Witch communes} with their patron at will as long as their favour is not zero.
Generally, a patron is helpful, but pursues their own goals and struggles to understand human ways outside of their domain. 
Favour determines how many times a witch can request an audience and expect an answer. 



\newcolumn

\subsection*{Uncovering a new spell} 

At first witch needs to convince the patron that they are ready by passing \texttt{Spirit} save. 

Witch loses amount of favour that is twice as much as the spell costs. It manifests as a small trinket which works as a puzzle, passing \texttt{Clarity} save is required to solve it.

After that, witch learns the new spell and can cast it as usual.


\end{multicols*}

